\section{Strong Induction}

\begin{frame}
\centering
\Huge\textbf{Strong Induction}
\end{frame}

% --------------------------------------------------
% Question 1
% --------------------------------------------------

\begin{frame}{Question 1}
A recursive algorithm can solve a problem of size $n$ by using a solution to any smaller positive problem size.

Suppose every integer $n\geq2$ can be reduced to either $n-1$ or $n-2$.

\textbf{Question:} Prove using strong induction that every positive integer can be reached from $1$ by repeatedly adding $1$ or $2$.
\end{frame}

\begin{frame}{Answer 1}
Let $P(n)$ be the statement:

\textit{$n$ can be reached from $1$ by repeatedly adding $1$ or $2$.}

\textbf{Base cases:}

$$
1
$$

is already reached.

Also,

$$
1+1=2.
$$

Thus, $P(1)$ and $P(2)$ are true.
\end{frame}

\begin{frame}{Answer 1}
\textbf{Strong inductive hypothesis:}

Assume that

$$
P(1),P(2),\ldots,P(k)
$$

are all true.

We must prove $P(k+1)$.

Since $k+1\geq2$, we can obtain it from the smaller integer $k$ by adding $1$:

$$
k+1=k+1.
$$

By the inductive hypothesis, $k$ can be reached from $1$.
\end{frame}

\begin{frame}{Answer 1}
Therefore, after reaching $k$, adding $1$ reaches $k+1$.

Hence $P(k+1)$ is true.

Thus, by strong induction,

$$
\boxed{\text{every positive integer can be reached from }1.}
$$

\end{frame}

% --------------------------------------------------
% Question 2
% --------------------------------------------------

\begin{frame}{Question 2}
A system represents an integer amount using coins of denominations $1$ and $2$.

\textbf{Question:} Prove using strong induction that every positive integer amount can be represented using these coins.
\end{frame}

\begin{frame}{Answer 2}
Let $P(n)$ denote:

\textit{$n$ can be represented using coins of denominations $1$ and $2$.}

\textbf{Base cases:}

$$
1=1
$$

and

$$
2=2.
$$

Therefore, $P(1)$ and $P(2)$ are true.
\end{frame}

\begin{frame}{Answer 2}
Assume that

$$
P(1),P(2),\ldots,P(k)
$$

are all true.

Consider $k+1$.

Since $k\geq2$,

$$
k+1=k+1.
$$

By the inductive hypothesis, $k$ can be represented using the available coins.

Adding one coin of denomination $1$ gives a representation of $k+1$.
\end{frame}

\begin{frame}{Answer 2}
Therefore, $P(k+1)$ is true.

Hence, by strong induction,

$$
\boxed{\text{every positive integer can be represented using coins of denominations }1\text{ and }2.}
$$

\end{frame}

% --------------------------------------------------
% Question 3
% --------------------------------------------------

\begin{frame}{Question 3}
A recursive program solves a problem of size $n$ by first solving problems of sizes smaller than $n$.

Suppose the program can solve the problem directly for sizes $1$ and $2$, and for every $n\geq3$, a solution for size $n-1$ is sufficient.

\textbf{Question:} Explain why strong induction is suitable for proving that the program works for every positive integer input size.
\end{frame}

\begin{frame}{Answer 3}
Let $P(n)$ denote:

\textit{The program correctly solves the problem for input size $n$.}

The base cases are

$$
P(1)
\quad\text{and}\quad
P(2).
$$

For the inductive step, assume

$$
P(1),P(2),\ldots,P(k)
$$

are all true.
\end{frame}

\begin{frame}{Answer 3}
For input size $k+1$, the program requires a correct solution for a smaller input size, specifically $k$.

Since

$$
P(k)
$$

is included in the strong inductive hypothesis, the required smaller problem is known to be solvable.

Therefore,

$$
P(k+1)
$$

is true.

Hence, strong induction proves correctness for every positive input size.
\end{frame}

% --------------------------------------------------
% Question 4
% --------------------------------------------------

\begin{frame}{Question 4}
Consider the Fibonacci sequence:

$$
F_1=1,\qquad F_2=1
$$

and

$$
F_n=F_{n-1}+F_{n-2}.
$$

\textbf{Question:} Prove using strong induction that

$$
F_n<2^n
$$

for every positive integer $n$.
\end{frame}

\begin{frame}{Answer 4}
Let

$$
P(n):\quad F_n<2^n.
$$

For $n=1$,

$$
F_1=1<2.
$$

For $n=2$,

$$
F_2=1<4.
$$

Thus, the base cases are true.
\end{frame}

\begin{frame}{Answer 4}
Assume strongly that

$$
F_i<2^i
$$

for every

$$
1\leq i\leq k.
$$

We prove the statement for $k+1$.

Using the recurrence,
$$
F_{k+1}=F_k+F_{k-1}.
$$

By the inductive hypothesis,
$$
F_k<2^k
$$

and
$$
F_{k-1}<2^{k-1}.
$$

\end{frame}

\begin{frame}{Answer 4}
Therefore,

$$
F_{k+1}
<
2^k+2^{k-1}
$$

$$
<2^k+2^k
$$
$$
=2^{k+1}.
$$

Hence,
$$
F_{k+1}<2^{k+1}.
$$

Therefore,
$$
\boxed{F_n<2^n}
$$

for every positive integer $n$.
\end{frame}

% --------------------------------------------------
% Question 5
% --------------------------------------------------

\begin{frame}{Question 5}
A file-compression algorithm represents files using blocks of size $2$ or $3$ units.

\textbf{Question:} Prove using strong induction that every integer $n\geq2$ can be represented as a sum of $2$s and $3$s, except $n=1$.
\end{frame}

\begin{frame}{Answer 5}
Let $P(n)$ denote:

\textit{$n$ can be expressed as a sum of $2$s and $3$s.}

The initial cases are

$$
2=2,
$$

$$
3=3,
$$

and
$$
4=2+2.
$$

Therefore,
$$
P(2),P(3),P(4)
$$

are true.
\end{frame}

\begin{frame}{Answer 5}
Assume strongly that

$$
P(2),P(3),\ldots,P(k)
$$

are all true.

Consider $k+1$.

For $k+1\geq5$,

$$
k+1=(k-1)+2.
$$

Since

$$
k-1\geq4,
$$

the strong inductive hypothesis guarantees that $k-1$ can be represented using $2$s and $3$s.
\end{frame}

\begin{frame}{Answer 5}
Adding one more $2$ gives a representation of $k+1$.

Therefore,

$$
P(k+1)
$$

is true.

Hence, by strong induction,

$$
\boxed{\text{every integer }n\geq2\text{ can be represented using }2\text{s and }3\text{s}.}
$$

\end{frame}
